\cxtitle{Commuting counterexamples to determinantal entropy concavity}

\begin{cxcredits}
\posedby{Russell Lyons \citeyearpar{Lyons2003DPM,Lyons2014DPP}}
\foundby{GPT Astra (Pro)}{2026-10-30}
\formalizedby{OpenAI Codex}
\auditedby{OpenAI Codex}
\contributedby{Suvrit Sra}
\end{cxcredits}

\begin{cxcontext}
Fix a finite coordinate set $E$. A positive contraction $Q$ is a self-adjoint
operator on the real or complex Hilbert space $\ell^2(E)$ satisfying
$0\preceq Q\preceq I$. Its determinantal law $\mathbf P^Q$ is defined by
$\mathbf P^Q[A\subseteq X]=\det Q_{A,A}$ for every $A\subseteq E$.
The entropy in the statement is the Shannon entropy of the entire occupied
set, with the coordinate basis fixed:
\[
 \operatorname{Ent}(Q)
 =-\sum_{S\subseteq E}\mathbf P^Q(\{S\})\log\mathbf P^Q(\{S\}),
 \qquad 0\log0=0.
\]
We use natural logarithms. Write $H(Q)=\operatorname{Ent}(Q)$ below. For the
coordinate projection $P_{S^c}$, the exact-set convention is
\[
 p_S(Q):=\mathbf P^Q(\{S\})
 =(-1)^{|S^c|}\det(Q-P_{S^c}).
\]
This is not the entropy of the particle count and not the sum of binary
entropies of the eigenvalues. The source explicitly permits complex Hilbert
spaces (Section~2.1 of \cite{Lyons2014DPP}), so complex Hermitian kernels are
within the stated domain; commutativity is a property of our witness, not a
hypothesis of the statement.
\end{cxcontext}

\begin{cxsource}{lyons-dpp-entropy-concavity}
R.~Lyons \citeyearpar{Lyons2014DPP}, Conjecture~2.6, eq.~(2.14); first posed
as Conjecture~9.2 \citep{Lyons2003DPM}
\end{cxsource}

\begin{cxstatement}{lyons-dpp-entropy-concavity}
For all positive contractions $Q_1$ and $Q_2$, we have
\[
 \operatorname{Ent}\bigl((Q_1+Q_2)/2\bigr)
 \ge \bigl(\operatorname{Ent}(Q_1)+\operatorname{Ent}(Q_2)\bigr)/2.
\]
\end{cxstatement}

\begin{cxsummary}{lyons-dpp-entropy-concavity}
The exact-set Shannon entropy of a determinantal law is midpoint concave in its positive-contraction kernel.
\end{cxsummary}

\begin{cxcertificate}{lyons-dpp-entropy-concavity}
Commuting $5\times5$ Hermitian kernels with rational spectra; rigorous rational logarithm bounds certify an entropy midpoint deficit exceeding $1/5{,}000{,}000{,}000$.
\end{cxcertificate}

\begin{cxrefutation}
The entropy conjecture of \cite{Lyons2014DPP} fails along a commuting affine
family. The matrices have an exact Fourier eigenbasis and rational spectra;
all squared minors needed for their occupation probabilities lie in
$\mathbb Q(\sqrt5)$.

\begin{theorem}[\statusfalse: determinantal entropy concavity]\label{thm:lyons-dpp-entropy-concavity}
Let $E=\{0,1,2,3,4\}$, let
\[
 \omega=e^{2\pi i/5},\qquad F_{jk}=5^{-1/2}\omega^{jk}
 \quad(0\le j,k\le4),
\]
and set
\[
 \lambda_+=\frac1{10000}(11,9,9991,9990,9989),\qquad
 \lambda_-=\frac1{10000}(9,11,9989,9990,9991),
\]
and $K_\pm=F\operatorname{diag}(\lambda_\pm)F^*$. These are commuting Hermitian strict positive contractions, but
\[
 H(K_+)=H(K_-),\qquad
 H(K_+)-H\left(\frac{K_++K_-}{2}\right)
 >\frac1{5{,}000{,}000{,}000}>0.
\]
In particular Conjecture~2.6 of \cite{Lyons2014DPP} is false as stated.
\end{theorem}
\begin{proof}
All displayed spectral entries are exact rational numbers. Because $F$ is
unitary, the two matrices commute, and their eigenvalues are strictly between
zero and one. Introduce the family
\[
 \lambda_\varepsilon(t)=
 (\varepsilon+t,\varepsilon-t,1-\varepsilon+t,
  1-\varepsilon,1-\varepsilon-t),\qquad
 K_\varepsilon(t)=F\operatorname{diag}(\lambda_\varepsilon(t))F^*.
\]
The theorem uses $\varepsilon=1/1000$ and $t=\pm1/10000$.
For every $0<\varepsilon<1/2$ and $|t|<\varepsilon$, the entire family
consists of commuting strict positive contractions.

For $J\subseteq E$ put
\[
 q_J(\lambda)=\prod_{j\in J}\lambda_j\prod_{j\notin J}(1-\lambda_j),
 \qquad M(S,J)=|\det F_{S,J}|^2.
\]
The exact-set probabilities satisfy
\begin{equation}\label{eq:lyons-dpp-spectral-probability}
 p_S(F\operatorname{diag}(\lambda)F^*)
 =\sum_{J:\,|J|=|S|}M(S,J)q_J(\lambda).
\end{equation}
Indeed the inclusion convention gives the probability generating polynomial
$\det(I-K+KZ)$, with $Z=\operatorname{diag}(z_i)$. Conjugating by $F$ and
expanding in the rows of $I-\operatorname{diag}(\lambda)$ and
$\operatorname{diag}(\lambda)F^*ZF$ gives
$\sum_Jq_J(\lambda)\det((F^*ZF)_{J,J})$. Cauchy--Binet expands the latter
minor as $\sum_{|S|=|J|}M(S,J)\prod_{i\in S}z_i$; comparison of coefficients
proves \eqref{eq:lyons-dpp-spectral-probability}.

Write $s=\sqrt5$. Squared Fourier minors of sizes zero and five equal one;
those of sizes one and four equal $1/5$. For two-element sets
$S=\{a,b\}$ and $J=\{c,d\}$,
\[
 M(S,J)=\frac{|\omega^{ac+bd}-\omega^{ad+bc}|^2}{25}
 =\begin{cases}
 (5-s)/50,&(b-a)(d-c)\equiv\pm1\pmod5,\\
 (5+s)/50,&(b-a)(d-c)\equiv\pm2\pmod5.
 \end{cases}
\]
The size-three values are the complementary size-two values, by the
complementary-minor identity for a unitary matrix. Thus every $p_S$ is an
explicit polynomial in $\varepsilon,t$ with coefficients in
$\mathbb Q(\sqrt5)$.

Let $R$ permute eigenmodes by $j\mapsto1-j\pmod5$. It fixes
$\lambda_\varepsilon(0)$ and reverses the direction $(1,-1,1,0,-1)$.
Since $FR=D\overline F$, where $D_{jj}=\omega^j$, we have
\[
 K_\varepsilon(-t)=D\overline{K_\varepsilon(t)}D^*.
\]
Complex conjugation and diagonal unitary conjugation preserve each exact-set
probability. Hence every $p_S(\varepsilon,t)$ is even in $t$, proving the
equality of the endpoint distributions and entropies.

It remains to establish the strictly positive finite entropy difference.
The accompanying \texttt{verify.py} evaluates
\eqref{eq:lyons-dpp-spectral-probability} with exact rational coefficients
and encloses all logarithms by rational bounds. It computes the eigenvalues
at both endpoints and checks strict admissibility, polynomial normalization,
and coefficientwise evenness. Its enclosure $[L,U]$ for
$H(K_+)-H((K_++K_-)/2)$ satisfies the exact assertions
\[
 L>\frac1{5{,}000{,}000{,}000}>0,
 \qquad
 L\le H(K_+)-H\left(\frac{K_++K_-}{2}\right)\le U.
\]
Here is the complete error-control method used by that finite computation.
With $D_0=10^{40}$ and $r=\lfloor\sqrt{5D_0^2}\rfloor$, integer square
root gives
\[
 r/D_0<\sqrt5<(r+1)/D_0.
\]
After exact power-of-two range reduction, logarithms of positive rational
endpoints reduce to arguments $1\le x\le2$. For
$z=(x-1)/(x+1)$ and $N=40$,
\[
 \log x=2\sum_{k=0}^{N-1}\frac{z^{2k+1}}{2k+1}+R_N,\qquad
 0\le R_N\le\frac{2z^{2N+1}}{(2N+1)(1-z^2)}.
\]
The series follows by integrating the geometric series; the tail bound
replaces each omitted denominator by $2N+1$. Algebraic arguments are
enclosed using the bounds on $\sqrt5$, and logarithms by monotonicity.
Interval multiplication uses the minimum and maximum of all four endpoint
products. Endpoints are rounded outward, using integer floor and ceiling,
to a grid of denominator $10^{32}$. These operations preserve containment
at each step. Every calculation and decisive inequality uses integers or
\texttt{fractions.Fraction}; no floating-point evidence is used. This proves
the stated finite inequality.
\end{proof}

There is also a complete analytic existence proof for the same family,
independent of executing the certificate. It explains the sign of the
curvature near a rank-three projection.

\begin{proposition}
For the family above,
\[
 \lim_{\varepsilon\downarrow0}
 \left.\frac{d^2}{dt^2}H(K_\varepsilon(t))\right|_{t=0}
 =2\log\frac{15}{44}
  +\frac2{\sqrt5}\log\frac{147+65\sqrt5}{22}>0.
\]
\end{proposition}
\begin{proof}
The kernel is circulant, so translating $S$ preserves its probability. For
the eight translation classes write
\[
 p_S(\varepsilon,0)=c_S\varepsilon^{m_S}(1+O(\varepsilon)),\qquad
 \partial_t^2p_S(\varepsilon,0)=d_S+O(\varepsilon).
\]
The Fourier-minor formula gives the following table; the last column is per
individual outcome, before multiplication by the class size.
\[
\begin{array}{c|r|r|c|c}
 S&\text{class size}&m_S&c_S&d_S\\\hline
 \varnothing&1&3&1&0\\
 \{0\}&5&2&3/5&-2/5\\
 \{0,1\}&5&1&(15-s)/50&(10-2s)/25\\
 \{0,2\}&5&1&(15+s)/50&(10+2s)/25\\
 \{0,1,2\}&5&0&(5-s)/50&(-10+6s)/25\\
 \{0,1,3\}&5&0&(5+s)/50&(-10-6s)/25\\
 \{0,1,2,3\}&5&1&2/5&4/5\\
 E&1&2&1&-2
\end{array}
\]
For completeness, the leading latent set at $\varepsilon=t=0$ is
$J_0=\{2,3,4\}$. The leading singleton contributions are its three
singletons, giving $3/5$. The leading pair contributions are
$\{2,3\},\{2,4\},\{3,4\}$, whose squared minors sum to
$(15\mp s)/50$. The leading triple contribution is $J_0$ itself. The two
leading four-point sets add either 0 or 1 to $J_0$, giving $2/5$.
The empty and full probabilities at $t=0$ are
$\varepsilon^3(1-\varepsilon)^2$ and
$\varepsilon^2(1-\varepsilon)^3$. These facts verify $m_S,c_S$.

To check $d_S$ directly, the only nonzero values of
$\partial_t^2q_J(\lambda_\varepsilon(t))$ at $\varepsilon=t=0$ are
\[
\begin{array}{c|r@{\qquad}c|r}
 J&q_J''&J&q_J''\\\hline
 \{3\}&-2&\{2,3\}&2\\
 \{3,4\}&2&\{0,2,3\}&2\\
 \{1,2,3\}&-2&\{0,3,4\}&-2\\
 \{1,3,4\}&2&\{2,3,4\}&-4\\
 \{0,2,3,4\}&2&\{1,2,3,4\}&2\\
 E&-2&&
\end{array}
\]
These are obtained by differentiating the five-factor products defining
$q_J$. Multiplication by the squared Fourier minors and summation gives
the displayed $d_S$ values. Polynomial dependence supplies all stated
error terms.

At every positive $\varepsilon$, all probabilities are positive, so entropy
is smooth. Evenness gives $p_S'(\varepsilon,0)=0$, and normalization gives
$\sum_Sp_S''=0$. Therefore
\[
 H_\varepsilon''(0)=-\sum_Sp_S''(\varepsilon,0)\log p_S(\varepsilon,0).
\]
The potential $\log\varepsilon$ terms cancel, because
$\sum_Sm_Sd_S=-4+4+4-4=0$; all errors are
$O(\varepsilon\log\varepsilon)$ or smaller and vanish. Put
$a=(5-s)/50$, $b=(5+s)/50$, $c=(15-s)/50$, $d=(15+s)/50$. The finite limit is
\begin{align*}
 C={}&2\log(3/5)-4\log(2/5)\\
 &-\tfrac25\bigl[(5-s)\log c+(5+s)\log d\bigr]
 +\tfrac25\bigl[(5-3s)\log a+(5+3s)\log b\bigr]\\
 ={}&2\log(15/44)+\tfrac2s\log\bigl((147+65s)/22\bigr),
\end{align*}
using $ab=1/125$, $cd=11/125$, and
$(c/d)(b/a)^3=(147+65s)/22$. Since
$15/44>1/3$, $(147+65\sqrt5)/22>13$, and $2/\sqrt5>8/9$,
\[
 C>\tfrac89\log13-2\log3
 =\tfrac29\log(28561/19683)>0.
\]
The second and third inequalities follow from $11/5<\sqrt5<9/4$.
Thus positive curvature occurs at interior kernels for every sufficiently
small positive $\varepsilon$. No boundary derivative of entropy is used.
\end{proof}

\begin{remark}[Scope]
The finite witness refutes the source's midpoint inequality itself, with
its original hypotheses and quantifiers intact. It also refutes full
concavity, even on a commuting affine family. It does not establish the
least possible dimension or settle the real-symmetric restriction.
It concerns finite exact-set Shannon entropy; it does not refute concavity
of particle-count entropy, spectral entropy, or any statement about entropy
rates of infinite stationary processes.

As a consequence, the auxiliary balanced beam-splitter inequality from the
discovery discussion also fails for this pair. For
\[
 W=\frac1{\sqrt2}\begin{pmatrix}I&I\\-I&I\end{pmatrix},\qquad
 Q=W(K_+\oplus K_-)W^*,\qquad \bar K=(K_++K_-)/2,
\]
both output coordinate marginals have kernel $\bar K$. Subadditivity yields
$H(Q)\le2H(\bar K)<H(K_+)+H(K_-)$, with output-minus-input gap strictly
less than $-1/2{,}500{,}000{,}000$. This is a corollary, not a second
provenance claim or an attribution of that auxiliary statement to Lyons.
\end{remark}
\end{cxrefutation}
